\subsection*{3. Number fields associated with a ``dessin d'enfant''}
The third section resumes the chronology of the surface questions and
constructs the passage from a finite map to an algebraic curve and its
arithmetic data. Ladegaillerie's isotopy work enters on pp.~9--10; the
``repères'', ``drapeaux'', or ``biarcs'' and the cartographic groups
$\underline C_2$ and $\underline C_2^+$ then give the dictionary for maps
and sets with an operator group. On p.~11 Grothendieck places the oriented
cartographic group in the fundamental-group description of the sphere
with three punctures, with the order-two condition at the middle
ramification value. The inverse image of $[0,1]$ recovers the map, and
the vertices, edge centers, and face centers become the points above
$0,1,\infty$ (p.~12); the footnote on that page also treats real curves
and unoriented maps. The rational function on p.~13 makes the Galois
action on coefficients explicit. He describes the discovery as a decisive
change in his mathematical center of interest on p.~12, and recounts
another change on pp.~14--15: the supposition he and Deligne had found
``dingue'' becomes the result announced by ``Bielyi''. The discussion then
names \emph{La Longue Marche à travers la théorie de Galois}, describes
the work of January--June 1981, and introduces the reconstruction questions
of ``géométrie algébrique anabélienne''. On p.~16 he records a change of
name from ``Théorie de Galois'' to ``théorie de Galois-Teichmüller'', while
returning to the faithful outer action already used in \S2. The p.~17
discussion carries these fundamental groups and their abelianizations
toward the representations associated with Jacobians. Its footnote (*)
contains a mathematical slip in the stated general description of the
quotient kernels as abelian products of Tate modules. The puncture-filling
morphism
$\pi_1(\mathbb P^1_{\overline{\mathbb Q}}\setminus\{0,1,\infty\})
\longrightarrow\pi_1(\mathbb P^1_{\overline{\mathbb Q}})=1$
has the entire free profinite group on two generators as kernel. The
assignment of those generators to $(12)$ and $(23)$ gives a continuous
surjection of this kernel onto $S_3$; their products are $(123)$ and
$(132)$ in opposite orders. A product of the abelian groups
$\hat{\mathbb Z}(1)$ cannot have that quotient, since the image of a
commuting pair remains commuting. This concrete morphism locates the
problem in the footnote's kernel description; its wording is preserved
in the readers. The closing paragraph on p.~18 returns to the role of
dreams and visions in mathematical work. Among subsequent texts one may
name Shabat and Voevodsky's
\href{https://www.math.ias.edu/vladimir/node/79}{\emph{Drawing Curves over Number Fields}}.
(\href{https://grothendieck.umontpellier.fr/web/viewer.html?file=../119.pdf\#page=12}{Copy 1, pp.~9--18, beginning at site 11};
\href{https://grothendieck.umontpellier.fr/web/viewer.html?file=../119.pdf\#page=20}{kernel description on p.~17, site 19}.)
% Complete links to notes (2), (3), and (4) after their scan checks.

